\scanpage{48}
\noindent\hfill atkin \mspage{8}

\[
\text{where}\quad\mathrm{P}=\handwritten{\prod_{i=1}^{\alpha}p_i^{\theta_i},}
\]
\noindent so that(since there are infinitely many primes) \uline{$c(n)$} is divisible by $13^{\alpha}$ in positive density of $n$ for every $\alpha\handwritten{\geq}1$. Hence in the final analysis the extensive computation of $c(13^{\handwritten{6}}n)$ becomes unnecessary\inserted{.} \deleted{\uline{for}} but without it we could not have found \handwritten{(3.9)}.

However in this connection there is a subtle difference in the corresponding work for \uline{$p(n)$},and generally for the Fourier coefficients of forms of half-integral dimension. In these cases there are results of M- and R-type,but they contain an unknown constant which is not a value of the coefficient,and one cannot prove the existence of ``accidents'' by the use of difference equations for a general prime $p\handwritten{\ne}13$. Thus in order to prove that \uline{$p(n)$} is divisible by $13^{\alpha}$ in positive density of \uline{$n$} it is necessary to compute $\alpha$ accidents,even though the \uline{general} M-result is proved \deleted{\uline{for all}} $\handwritten{\unsure{\alpha|u}}$.\margin{\handwritten{\tikz[baseline=0pt,x=.3em,y=.4ex,line cap=round]{\draw (.75,3.9)..controls(.8,3.15)and(.1,3.3)..(.2,2.6)..controls(.3,2.05)and(1,2.4)..(.72,2.8)..controls(.4,3.1)and(.25,1.8)..(.7,1.55)..controls(1.05,1.25)and(.45,.9)..(.28,1.3);\draw(.05,-.6)--(.83,4.1);}}} I have in fact found 13 such accidents,so that \uline{$p(n)$} is divisible by $13^{13}$ in positive density of $n$. (The second ``13'' is of course an irrelevant accident!)

\noindent \deleted{\unsure{1|2}}\handwritten{4}. \uline{The} \uline{cases} $q=17,19$,\uline{and} $23$. Concurrently with the work for $\uline{q}=13$ described in the last section,the methods of ATKIN(1967) were used to obtain suitable systems of modular functions for the primes $\uline{q}=17,19,23,29,31$,and $37$,with a view to obtaining \deleted{\unsure{24|27|p)}} \margin{\handwritten{\tikz[baseline=0pt,x=.35em,y=.4ex,line cap=round]{\draw (.45,3.8)..controls(.42,3.1)and(-.05,3.15)..(.02,2.5)..controls(.08,2)and(.64,2.2)..(.52,2.7)..controls(.44,2.95)and(.03,1.8)..(.4,1.5)..controls(.75,1.15)and(.25,.9)..(.1,1.3);\draw(.15,-.5)--(1.45,3.8);}}}R-results. As soon as the M-results for 13 appeared,it required no additional personal effort(and in fact very little machine time)to test the situation on these primes with regard to M-results. I found that the conjectures \handwritten{(3.9)} and \handwritten{(3.10)} held for $\uline{q}=17,19$,and $23$,word for word but replacing 13 by \uline{$q$}. This provides a further illustration of the value of \handwritten{a} computer: with the necessary subroutines already written one can multiply instances without the loss of momentum which hand calculation would involve. There was in this case no particular reason to expect ``the same again''; the type of modular equation for $\uline{\unsure{q|c}}=13$ does not exist for $\uline{\unsure{q|c}}>13$\unsure{:|.}

% archive sheet number: 47
